% ============================================================
% CAS/Python ενότητες — Κεφάλαιο 12: Εφαρμογές Ολοκληρωτικού
%
% QR code: qr_3c203fbbafb7.png  (→ latex/qrcodes/)
% URL: https://nikmatz.github.io/ELADIC_GR/code/ch12/
% ============================================================

\casactivbox%
  {Maxima: Εμβαδόν, Όγκος Περιστροφής, Μήκος Τόξου}%
  {%
    \label{cas:ch12_integral_apps}
    \begin{enumerate}[label=(\alph*),itemsep=2pt]
      \item \textbf{Εμβαδόν}: υπολογίστε με \texttt{integrate} και
        \texttt{romberg} το $\int_{-1}^{1}|x^3-x|\,dx$ και το
        εμβαδόν μεταξύ $y=-x^2+4$ και $y=x^2-2x$.
        Βρείτε τομές με \texttt{solve}.
      \item \textbf{Όγκος περιστροφής} — μέθοδος δίσκων
        ($V=\pi\int_a^b[f(x)]^2\,dx$): υπολογίστε για
        $y=\sqrt{x}$ σε $[0,4]$ και $y=\sin x$ σε $[0,\pi]$.
        Δακτύλιοι: $y=\sqrt{x}$ vs $y=x$ γύρω από τον $x$-άξονα.
        Κελύφη: $y=x^3$ γύρω από τον $y$-άξονα, $[0,2]$.
      \item \textbf{Μήκος τόξου} ($L=\int_a^b\!\sqrt{1+[f'(x)]^2}\,dx$):
        για $y=x^{3/2}$ σε $[0,4]$ και $y=\ln(\cos x)$ σε $[0,\pi/3]$.
        Χρησιμοποιήστε \texttt{diff} και \texttt{trigsimp}.
      \item \textbf{Εμβαδόν επιφάνειας} ($S=2\pi\!\int f\sqrt{1+f'^2}\,dx$):
        για τον κώνο $y=x$, $[0,2]$. Επαληθεύστε με τον τύπο
        $S=\pi r l$ ($r=2$, $l=2\sqrt{2}$).
      \item \textbf{Φυσικές εφαρμογές}: έργο ελατηρίου $F=200x$,
        $x\in[0,0.1]$· κέντρο βάρους $y=x^2$, $[0,3]$.
    \end{enumerate}
    \smallskip
    \textbf{Εντολές Maxima:}
    \texttt{integrate(f,x,a,b)}, \texttt{romberg(f,x,a,b)},
    \texttt{diff(f,x)}, \texttt{sqrt()}, \texttt{solve(f=g,x)}
  }%
  {https://nikmatz.github.io/ELADIC_GR/code/ch12/}%
  {qr_3c203fbbafb7.png}

\subsection*{Δραστηριότητα Python}

\begin{pybox}{Python: Εφαρμογές Ολοκληρωτικού — SymPy, SciPy, 3D}%
             {https://nikmatz.github.io/ELADIC_GR/code/ch12/}%
             {qr_3c203fbbafb7.png}
\label{py:ch12_integral_apps}
\begin{enumerate}[label=(\alph*),itemsep=2pt]
  \item Υπολογίστε εμβαδόν μεταξύ καμπυλών με
    \texttt{sympy.integrate} και σχεδιάστε τη σκιασμένη
    περιοχή με \texttt{matplotlib.fill\_between}.
  \item Όγκος περιστροφής: εφαρμόστε και τις τρεις
    μεθόδους (δίσκων, δακτυλίων, κελυφών) και επαληθεύστε
    ότι δίνουν το ίδιο αποτέλεσμα. Οπτικοποιήστε το στερεό
    3D με \texttt{mpl\_toolkits.mplot3d}.
  \item Μήκος τόξου: υπολογίστε με \texttt{sympy} και
    \texttt{scipy.integrate.quad}. Σχεδιάστε σωρευτικό
    μήκος τόξου $L(t)$ παράλληλα με την καμπύλη.
  \item Φυσικές εφαρμογές: υπολογίστε έργο ελατηρίου,
    κέντρο βάρους και υδροστατική δύναμη σε ορθογωνική
    δεξαμενή ($4\text{ m}\times 3\text{ m}$, $\rho g=9810$ N/m³).
\end{enumerate}
\medskip
\textbf{Βιβλιοθήκες / Εντολές:}
\texttt{sympy.integrate()}, \texttt{sympy.diff()},
\texttt{scipy.integrate.quad()},
\texttt{mpl\_toolkits.mplot3d}, \texttt{fill\_between()}
\end{pybox}
